\section{Introduction}

River-mouth morphodynamics involve coupled water motion, suspended sediment transport, deposition, and bed change. A full field-scale model must handle nonuniform beds, open boundaries, tides, wetting and drying, erosion, turbulence closures, and calibration data. This paper deliberately does not attempt that full problem. Instead, it isolates a restricted numerical baseline for which several discrete properties can be stated and proved without relying on observational data or physical validation.

The long-term motivation is an idealized two-dimensional river-mouth setting in which a channel discharges into a receiving basin. At this stage, however, the mathematical result is only for a periodic rectangular domain with a flat fixed hydrodynamic bed. Suspended sediment is represented by the conservative variable
\[
  s = h c,
\]
where \(h\) is water depth and \(c\) is a dimensionless volumetric concentration. Deposition removes suspended sediment locally and stores the same volume in a separate variable \(\beta\). The storage variable is only a sediment bookkeeping variable in the theorem; it is not used as a hydrodynamic bed elevation.

The contribution of this paper is therefore modest and specific. We define a first-order finite-volume scheme using a Rusanov shallow-water update, a first-order upwind sediment advection update, and an exact deposition step. Under periodic boundaries, strict wetness, and explicit CFL restrictions, we prove one-step water-depth positivity, nonnegative suspended sediment and concentration, exact water-mass conservation, and exact conservation of total suspended-plus-stored sediment volume in exact arithmetic.

The paper does not claim novelty, convergence of the full morphodynamic model, well-balancedness, energy stability, or validation for any real river. Existing work already contains important results on positivity-preserving and well-balanced shallow-water schemes, scalar transport, and Exner-type sediment mass balance \cite{audusse2004,kurganov2007,paola2005,karjoun2020}. The present baseline is best viewed as a transparent proof and implementation target from which later extensions can be assessed.
